\subsection{Quarterly support and estimand}
\label{sec:si-quarterly-dynamics}

The quarterly observation frame contains all 334 cities and 32 natural quarters in
2016--2023.  Eligibility requires finite current $Q\geq0$, $E>0$ and
$W>0$, and finite positive $N$ in quarters $t,t-1,t-2,t-3$.  Lag values are
joined by city identifier and integer calendar-quarter ordinal; removing an
unobserved quarter therefore cannot create a pair from two nonadjacent
observations.  Lagged activity may be zero or missing because the lag criterion
concerns facilities.  Current source-coded-zero activity is retained, whereas
absent compatible weeks remain missing.  All-zero fixed-effect groups,
singleton groups and possible statistical separation are recorded separately
from missing observations.

The model is for $Q$ with $\ln W$ as a fixed offset.  It estimates a conditional
arithmetic weekly mean,
\[
 \mathbb E[Q\mid X,W]=W\exp(Xb),
\]
so the data need not be integer counts.  Exposure affects the score and Hessian
even without an additional frequency-weight argument.  Complete calendar-
quarter fixed effects absorb common time shocks, including common seasonality.
No quarterly weather field is observed in this source panel; annual weather or
background columns are not copied into quarterly controls.

The static, one-quarter and full three-lag specifications retain the
intersection of their estimable supports.  Each uses city and calendar-quarter
effects and the same observations; the one-quarter specification does not gain
rows from its shorter lag.  The pandemic comparison starts within this support,
restricts years to 2022 or earlier, requires complete current source-date
 exposure, and compares specifications without and with the response index on
 the same support with province clustering.  The latter adds $\mathrm{COVID}/10$.
A source-designated pre-pandemic structural zero is distinct from an unobserved
index, and missing 2023 values are never filled.  The paired comparison is
therefore between two specifications on a common subset, not between that
subset and the full-period estimate.

\subsection{Equivalent parameterization and finite-lag contrasts}
For the full three-lag specification, the level equation in the main text can
be written exactly as
\begin{align}
 \log\{\mathbb E[Q_{it}\mid X,W]/W_{it}\}
 ={}&\alpha_i+\tau_t+\theta x_{i,t-3}
 +\delta_0(x_{it}-x_{i,t-1})\\
 &+\delta_1(x_{i,t-1}-x_{i,t-2})
 +\delta_2(x_{i,t-2}-x_{i,t-3}).
\end{align}
The coefficient mapping is
$\beta_0=\delta_0$, $\beta_1=\delta_1-\delta_0$,
$\beta_2=\delta_2-\delta_1$ and $\beta_3=\theta-\delta_2$; hence
$\theta=\sum_{j=0}^3\beta_j$.  If $T$ maps fitted coordinates to level
coefficients, then $b_{\mathrm{level}}=Tb_{\mathrm{fit}}$ and
$V_{\mathrm{level}}=TV_{\mathrm{fit}}T^{\top}$.  The base-level term is
essential: a regression containing only growth terms would not estimate the
facility-level response.  Four consecutive facility observations are required
for this specification.

The stipulated comparison raises facilities by a common proportion from
quarter zero onward while holding the fixed-effect and other covariate paths
fixed.  At horizon $h$, the conditional weekly-mean elasticity is
$c_h=\sum_{j=0}^h\beta_j$.  It is a quarter-specific level response, not an
elasticity of charging accumulated across quarters.  For a 10\% facility
difference, the activity contrast is
$100[1.1^{c_h}-1]$ and its interval transforms the endpoints of the elasticity
interval.  The corresponding intensity contrast is
$100[1.1^{c_h-1}-1]$.  Neither calculation extrapolates an empirically estimated
response beyond the fixed lag horizon.

Let $a_h$ select the first $h+1$ level coefficients.  Then
$\widehat c_h=a_h^{\top}\widehat b$ and
$\mathrm{se}(\widehat c_h)=\sqrt{a_h^{\top}\widehat V a_h}$.  The extra
follow-through contrast is $c_3-c_0=\beta_1+\beta_2+\beta_3$, with its own
covariance-based standard error.  The two primary two-sided tests are
$c_3=1$ and $c_3-c_0=0$.  Pointwise 95\% intervals for the remaining path
points are not simultaneous bands.  The one-quarter specification concerns one
subsequent quarter, whereas the full distributed-lag specification concerns
three; the static specification has no lag.

\subsection{Cluster covariance and fit diagnostics}
For the full design $X$, fitted means $\widehat\mu$ and cluster score sums
$s_g=\sum_{i\in g}X_i^{\top}(Q_i-\widehat\mu_i)$, the covariance is
\begin{equation}
 \widehat V=\frac{G}{G-1}\frac{n-1}{n-r}
 (X^{\top}\operatorname{diag}(\widehat\mu)X)^{-1}
 \left(\sum_{g=1}^{G}s_gs_g^{\top}\right)
 (X^{\top}\operatorname{diag}(\widehat\mu)X)^{-1},
\end{equation}
where $r$ is the full design rank, including fixed effects and substantive
regressors.  Main models cluster by city; the pandemic pair clusters by
province.  Contrasts use two-sided $t(G-1)$ critical values with the relevant
cluster count.  Full coefficients and covariance blocks are retained,
including transformations from numerical scaling or equivalent coordinates.

The support, fixed-effect admissibility and separation checks below determine
which observations and covariance calculations enter these estimands.

The quarterly estimates use Python \texttt{statsmodels.GLM} with the Poisson
family and iteratively reweighted least squares, followed by explicit
Poisson-score/Hessian cluster covariance calculations.  City all-zero groups
and singletons are removed before fitting.  A linear program checks whether a
design direction leaves positive-outcome predictors unchanged while weakly
reducing zero-outcome predictors, with strict reduction for at least one zero;
this is the separation criterion discussed by
\citet{correia2026separation}.  No strict separation direction is detected in
any of the five designs.  The complete covariance is transformed from the
fitted base-level/difference parameterization to the original-level
coefficients.

All specifications converge.  The common main support contains 1,301 positive
and 241 source-coded-zero outcomes: 1,542 rows in 111 cities over 28 calendar
quarters from 2017Q1 through 2023Q4, with 2--28 retained quarters per city and
a median of 11.  Every retained row has $W/W^{\mathrm{cal}}=1$; the minimum,
quartiles and maximum are all one.  Complete coverage was not imposed as an
additional sample threshold.  Main-model inference uses $t(110)$.

Before fixed-effect support exclusions, the four-quarter facility support has
2,229 rows in 273 cities, including 1,309 positive and 920 source-coded-zero
outcomes.  The main support removes 155 all-zero city groups (679 rows) and 7
singleton city groups (7 rows); one additional row fails the city/quarter
fixed-effect support condition.  The result is the 1,542 city-quarter support
in 111 cities.

Before fixed-effect support exclusions, the COVID support has 1,107 rows in 104
cities.  Fixed-effect support exclusions remove 17 rows: 13 rows meet the
city-all-zero condition and 5 meet the city-singleton condition, with one row
meeting both conditions.  Equivalently, the mutually exclusive count is 13
all-zero rows plus 4 additional singleton-only rows.  The resulting paired
support has 1,090 rows in 94 cities over 24
quarters and 25 province clusters; its inference uses $t(24)$.  No 2023 COVID
exposure is filled.  These exclusions define the reported estimands and are
separate from missing compatible weeks.

